% ============================================================
% Weekly Meeting Report — Week 22
% Topic: Stage 3 Launch, Scope Review, Thursday Supervisor Meeting
% ============================================================

\section*{Important - Feasibility - Timeline}
\addcontentsline{toc}{section}{Important - Feasibility - Timeline}

\begin{tabular}{ll}
    \textbf{Date}       & 5~June~2026 \\
    \textbf{Student}    & Khac Duc Giang Nguyen \\
    \textbf{Supervisor} & Dr.\ Seyed Sahand Mohammadi Ziabari \\
    \textbf{Phase}      & Stage~3 Running; Scope Review \& Thesis Alignment \\
\end{tabular}

\vspace{1em}

% ------------------------------------------------------------
\subsection*{1. Overview}

This report covers the period 2--5~June~2026 (Week~22). Stage~3
training was launched on 2~June~2026 with three parallel conditions
(naive, random stratified, herding). This report also documents a
formal scope review of the thesis against the original design
document (submitted 8~March~2026), identifying what has been
delivered, what can still be added, and what is out of scope for
the June~26 submission.

% ------------------------------------------------------------
\subsection*{Compute Budget — Exhausted, No Extension}

\begin{table}[h]
\centering
\small
\begin{tabular}{ll}
\hline
Initial budget      & 120,000 SBU \\
Used to date        & $\approx$113,140 SBU \\
\textbf{Remaining}  & \textbf{6,860 SBU} \\
\hline
\end{tabular}
\end{table}

\noindent No additional compute budget will be available before the
June~26 deadline. The remaining 6,860\,SBU is reserved for
post-training analysis jobs (evaluation scripts, figure generation)
running on the cheaper A100 partition ($\approx$128\,SBU/h). No
further training runs are possible.

\subsubsection*{Stage 3 outcome}

The herding condition was submitted four times (jobs 23425651,
23442877, 23443056, 23449442). Three successive bugs were identified
and fixed across these attempts:

\begin{enumerate}
    \item \textbf{Job~23425651}: PyTorch autograd in-place conflict
    under DDP — fixed with \texttt{model.no\_sync()}.
    \item \textbf{Job~23442877}: \texttt{TypeError} in logging f-string
    (tensor formatted without \texttt{.item()}).
    \item \textbf{Job~23443056}: Silent data bug — \texttt{ReplayDataset}
    looked for \texttt{meta['img\_path']} which does not exist in
    \texttt{AntiUAVRGBTDataset} metas. Every replay image loaded as
    a black zero tensor. Fixed by reconstructing frames from
    \texttt{meta['seq\_dir']} and \texttt{meta['frame']}.
    \item \textbf{Job~23449442}: Submitted with all fixes applied.
    Cancelled after the budget situation was assessed — insufficient
    SBU to run to completion.
\end{enumerate}

\noindent The code is correct. The experiment simply ran out of time
and money. \textbf{Stage~3 herding will not produce results for this
submission.}

% ------------------------------------------------------------
\subsection*{2. Current Experimental Status}

\begin{table}[h]
\centering
\caption{Full experimental status as of 5~June~2026.}
\label{tab:status}
\small
\begin{tabular}{p{0.38\linewidth}p{0.12\linewidth}p{0.42\linewidth}}
\hline
\textbf{Component} & \textbf{Status} & \textbf{Notes} \\
\hline
\multicolumn{3}{l}{\textit{Training pipeline}} \\
Stage~1 (Anti-UAV-RGBT)         & \checkmark Done  & mAP@0.5\,=\,0.6725, T1 baseline locked \\
Stage~2 KD fine-tuning (Anti-UAV410), 3 seeds & \checkmark Done  & 0.428\,$\pm$\,0.002; FM\,=\,$-$0.033\,$\pm$\,0.004 \\
Stage~3 naive baseline          & \checkmark Done  & FM\_naive\,=\,$-$0.6045, 18 epochs \\
Stage~3 random stratified (ablation) & Cancelled & Budget exhausted \\
Stage~3 herding (contribution)  & Cancelled & Budget exhausted (4 crash-fix cycles) \\
\hline
\multicolumn{3}{l}{\textit{Analysis \& evaluation}} \\
Stage~1 full analysis           & \checkmark Done  & Scale bins corrected; mAP\,=\,0.6728 \\
Stage~2 multi-seed CI plots     & \checkmark Done  & $\pm$\,std across 3 seeds \\
Stage~2 scale-stratified eval   & \checkmark Done  & Normal: 0.681, Tiny: 0.001 \\
Stage~2 parameter drift (S1$\to$S2) & \checkmark Done & mean cosine sim\,=\,0.91; spatial attention drifts most \\
Stage~3 eval (post-training)    & Pending  & Run \texttt{eval\_full\_analysis.py} after jobs finish \\
\hline
\multicolumn{3}{l}{\textit{Data \& infrastructure}} \\
Replay buffers (herding + random) & \checkmark Done  & 300 exemplars each, Stage~2 features \\
Dataset loaders (all 3 datasets) & \checkmark Done  & RGBT, UAV410, CST \\
\hline
\end{tabular}
\end{table}

% ------------------------------------------------------------
\subsection*{3. Scope Review Against Original Design}

The original thesis design proposed a tri-modular pipeline. This
section documents the actual scope delivered and the justification
for each deviation.

\subsubsection*{3.1 What Has Been Fully Delivered}

\begin{enumerate}
    \item \textbf{Sequential continual learning pipeline (RQ2, H1).}
    Three-stage curriculum: Anti-UAV-RGBT $\to$ Anti-UAV410 $\to$
    CST Anti-UAV. Forgetting Measure tracked at every epoch.
    This is the central scientific contribution.

    \item \textbf{Scale-Stratified Herding replay buffer.}
    iCaRL-style greedy herding in the Stage~2 feature space,
    stratified equally across four UAV size bins. Both buffers
    (herding and random) built and code-verified. Experimental
    validation was blocked by compute budget exhaustion.

    \item \textbf{Multi-seed confidence intervals.}
    Stage~2 results validated across 3 seeds
    (T2 mAP\,=\,0.428\,$\pm$\,0.002). Stage~3 naive ran 18 epochs
    on a single seed; no further Stage~3 runs are possible.

    \item \textbf{Knowledge Distillation (teacher-student) in Stage~2.}
    Stage~1 model acts as teacher; Stage~2 student trained with
    combined detection + KD loss. KD/det ratio rises from 1.36$\times$
    to 2.60$\times$ — confirms KD constraint is active.
\end{enumerate}

\subsubsection*{3.2 Three-Condition Ablation — Designed, Not Executed}

The experimental design included three parallel Stage~3 conditions:
naive baseline, random stratified replay, and herding replay. The
random stratified condition was specifically added (beyond the original
design) to isolate the contribution of strata balance from the herding
selection algorithm — a methodological strengthening of H1. All three
conditions were implemented and the code was debugged across four
submission cycles. However, only the naive condition produced results;
random stratified and herding were cancelled when the compute budget
was exhausted. The three-condition design is presented in the
methodology chapter as the intended ablation; results are reported for
naive only, with herding and random stratified as future work.

\subsubsection*{3.3 RQ1 — Delivered Differently, Still Valid}

RQ1 as proposed asked whether a Teacher-Student framework could bridge
the \textbf{RGB$\to$TIR modality gap} without manual annotations,
using the RGB stream from Anti-UAV300 as teacher and TIR as student.

\noindent What was actually implemented is a related but distinct
contribution: \textbf{same-modality domain-incremental Knowledge
Distillation}. The Stage~1 TIR model acts as teacher during Stage~2,
constraining the student from forgetting T1 representations while
adapting to Anti-UAV410. This is TIR$\to$TIR, not RGB$\to$TIR, but
the teacher-student mechanism and the anti-forgetting objective are
both present and measurable.

\noindent \textbf{Justification for why this still addresses the spirit
of RQ1:}

\begin{enumerate}
    \item The original motivation for RQ1 was to preserve source-domain
    knowledge while adapting to a new domain — \textit{without} requiring
    full re-annotation of the target domain. This is exactly what the
    KD loss achieves: the Stage~2 model adapts to Anti-UAV410 using only
    Anti-UAV410 labels, while the KD loss implicitly enforces retention
    of Anti-UAV-RGBT representations through soft supervision from the
    Stage~1 model.

    \item The cross-modality gap (RGB$\to$TIR) was not implemented because
    Anti-UAV-RGBT was used as the Stage~1 training dataset in TIR mode
    from the start. The dataset does contain paired RGB streams, but
    the thesis scope shifted to domain-incremental adaptation (dataset
    shift, same modality) rather than modality adaptation. This is a
    narrower but cleaner experimental design that directly isolates the
    scale-shift forgetting effect tested in H1.

    \item The parameter drift analysis provides direct empirical evidence
    that the KD mechanism worked as intended: mean cosine similarity of
    0.91 between Stage~1 and Stage~2 weights confirms the teacher
    successfully constrained the student. The spatial attention module
    (model.4) shows the highest drift ($>4\times$ relative L2), indicating
    the model adapted its attention patterns for the new domain while
    preserving backbone representations — the desired stability-plasticity
    balance.
\end{enumerate}

\noindent \textbf{Proposed reframing of RQ1 for the thesis:}
\begin{quote}
\textit{``To what extent does Knowledge Distillation from a
source-domain model mitigate catastrophic forgetting during
sequential domain adaptation, and how is this preservation
distributed across network layers?''}
\end{quote}

\noindent This reframing is fully answered by the existing results:
FM\,=\,$-$0.033\,$\pm$\,0.004 and the parameter drift analysis
together provide a quantitative and mechanistic answer. The
cross-modal UDA component is repositioned as future work, with
the note that Anti-UAV-RGBT's paired RGB streams make it
directly feasible in a follow-up study.

\subsubsection*{3.4 What Can Be Added Before Submission (Short Time)}

These require no retraining — only evaluation scripts.

\begin{enumerate}
    \item \textbf{Tracking metrics: SR, PR on CST sequences.}
    Write an evaluation script that runs \texttt{best.pt} on CST
    val sequences and computes Success Rate and Precision using
    IoU-Hungarian matching. Addresses RQ3 partially (without the
    evidential head).

    \item \textbf{ARD100 out-of-distribution evaluation.}
    Evaluate Stage~1 and Stage~2 \texttt{best.pt} on ARD100
    sequences as a stress test. No training required. Provides
    the ``ego-motion robustness'' data point from the original
    ablation table.

    \item \textbf{Per-stratum forgetting analysis (already in code).}
    \texttt{train\_stage3.py} already logs mAP@0.5 per size stratum
    (tiny/small/normal/large) on the T1 val set at every epoch.
    After Stage~3 completes, these curves show exactly which strata
    are forgotten by naive training and preserved by herding.
\end{enumerate}

\subsubsection*{3.5 Feasible But Tight}

\begin{enumerate}
    \item \textbf{RGB stream as \texttt{img2} for Stage~1 (retrain S1+S2+S3).}
    Anti-UAV-RGBT (= Anti-UAV300) contains paired RGB and TIR streams.
    Using the RGB frame as \texttt{img2} would activate YOLOMG's
    dual-input branch with real appearance data rather than zeros.
    This requires retraining Stage~1 ($\sim$14\,h), Stage~2
    ($\sim$30\,h), and Stage~3 ($\sim$30\,h) — approximately
    3.5 days of compute. Feasible before June~26 if decided
    immediately after this meeting.
    \textbf{Decision required from supervisor.}
\end{enumerate}

\subsubsection*{3.6 Out of Scope — Justified}

\begin{table}[h]
\centering
\caption{Components out of scope with justifications.}
\label{tab:out_of_scope}
\small
\begin{tabular}{p{0.25\linewidth}p{0.68\linewidth}}
\hline
\textbf{Component} & \textbf{Justification} \\
\hline
GMC + motion difference masks
& \textbf{Anti-UAV410 is an image sequence with no continuous video}
  (individual frames), making motion masks technically impossible for
  Stage~2. Since consistent \texttt{img2} across all stages is required,
  zeros is the only pipeline-consistent choice. ARD100 is the intended
  use case for GMC but falls outside this thesis's dataset curriculum. \\[4pt]

Evidential Head + KL regularisation
& Requires architectural modification to YOLOMG and retraining from
  scratch. Given Stage~3 is already running and submission is
  June~26, this is not feasible. Positioned as future work. \\[4pt]

Full RGB$\to$TIR Teacher-Student UDA (original RQ1)
& The proposed UDA used RGB from Anti-UAV300 to supervise TIR
  feature extraction. In the implemented pipeline, Stage~2 performs
  TIR$\to$TIR KD (same modality, different dataset). The cross-modality
  gap is addressed implicitly by the dataset curriculum, not explicitly
  by a UDA loss. RQ1 has been reframed accordingly (Section~3.2). \\[4pt]

ARD100 as training data
& ARD100 was designated for motion mask generation and ego-motion
  stress testing. Without the motion branch, it has no training role.
  It can still be used for out-of-distribution evaluation (Section~3.3). \\[4pt]

ECE calibration
& Requires the evidential head, which is out of scope. \\[4pt]

IDSW reduction via bimodal fusion
& Requires both the evidential head and real motion masks.
  Out of scope. \\
\hline
\end{tabular}
\end{table}

% ------------------------------------------------------------
\subsection*{4. Motion Branch Justification (Detailed)}

The motion branch (\texttt{img2} channel in YOLOMG) is set to zeros
across all three stages. This is a deliberate, principled decision
for the following dataset-specific reasons:

\begin{enumerate}
    \item \textbf{Anti-UAV-RGBT (Stage~1):} Video sequences with
    reactive camera movement. The camera operator manually tracks
    the drone, causing sporadic pan/tilt events. Without GMC,
    inter-frame differences during pan events are dominated by
    background noise. Since stable and pan frames cannot be
    reliably distinguished at training time, feeding raw motion
    masks would introduce inconsistent noisy signal — worse than
    zeros.

    \item \textbf{Anti-UAV410 (Stage~2):} An image sequence dataset.
    Frames within a sequence are ordered but the dataset does not
    guarantee continuous video coverage. Motion masks are technically
    inadvisable. More critically, \textbf{this is the hard blocker}
    for pipeline-wide motion mask adoption: even if motion masks were
    computed for Stage~1 and Stage~3, they cannot be computed for
    Stage~2, making consistent dual-input training impossible.

    \item \textbf{CST Anti-UAV (Stage~3):} Video sequences with a
    \textbf{static camera}. Motion masks are trivially computable
    (simple frame difference, no GMC required). However, enabling
    the motion branch only in Stage~3 while it was disabled in
    Stage~1 and Stage~2 would introduce a new input modality that
    the model has never been trained to process, injecting noise
    rather than signal.
\end{enumerate}

\noindent The conclusion is that consistent \texttt{img2}\,=\,0
across all stages is the only scientifically sound choice given
the dataset constraints. The motion branch is reserved for a
future pipeline where all training stages use datasets with
suitable temporal sequences (e.g., replacing Anti-UAV410 with
a video-native thermal dataset).

% ------------------------------------------------------------
\subsection*{5. Key Numbers for the Thesis}

\begin{table}[h]
\centering
\caption{Headline results to date (Stage~3 pending).}
\label{tab:headline}
\small
\begin{tabular}{llr}
\hline
\textbf{Stage} & \textbf{Metric} & \textbf{Value} \\
\hline
Stage~1 & T1 mAP@0.5 (ceiling)              & 0.6725 \\
Stage~1 & Precision / Recall / F1            & 0.636 / 0.661 / 0.648 \\
Stage~1 & Normal-stratum mAP@0.5             & 0.719 \\
Stage~1 & Tiny-stratum mAP@0.5              & 0.009 \\
\hline
Stage~2 & T2 mAP@0.5 (mean $\pm$ std, n=3)  & 0.428 $\pm$ 0.002 \\
Stage~2 & FM at peak T2 epoch (corrected)    & $-$0.033 $\pm$ 0.004 \\
Stage~2 & Normal-stratum mAP@0.5             & 0.681 \\
Stage~2 & Tiny-stratum mAP@0.5              & 0.001 \\
Stage~2 & KD/det loss ratio (final)          & $\approx$2.60$\times$ \\
Stage~2 & Mean cosine sim (S1$\to$S2)        & 0.917 \\
\hline
Stage~3 & FM\_naive (fm\_abs, at peak T3 ep.) & $-$0.6045 \\
Stage~3 & FM\_naive (fm\_stage3)             & $-$0.5724 \\
Stage~3 & Overall T1 mAP@0.5 (post-training eval) & 0.0761 \\
Stage~3 & Normal-stratum T1 mAP@0.5         & 0.0878 (from 0.719) \\
Stage~3 & Large-stratum T1 mAP@0.5          & 0.0000 \\
Stage~3 & Small-stratum T1 mAP@0.5          & 0.0670 \\
Stage~3 & Tiny-stratum T1 mAP@0.5           & 0.0004 \\
Stage~3 & FM\_random / FM\_herding           & N/A — not run \\
\hline
\end{tabular}
\end{table}

% ------------------------------------------------------------
\subsection*{6. Research Question Status}

\begin{table}[h]
\centering
\caption{Research question status at submission.}
\small
\begin{tabular}{p{0.08\linewidth}p{0.22\linewidth}p{0.14\linewidth}p{0.46\linewidth}}
\hline
\textbf{RQ} & \textbf{Question (short)} & \textbf{Status} & \textbf{Notes} \\
\hline
H1 &
Catastrophic forgetting is measurable across sequential tasks? &
\textbf{Confirmed} &
FM\,=\,$-$0.605 at peak T3 epoch. Large-stratum T1 mAP\,$\to$\,0.000 by epoch~1. \\[4pt]

RQ1 &
Does KD preserve T1 performance during T2 domain shift? &
\textbf{Fully answered} &
Stage~2 FM\,=\,$-$0.033\,$\pm$\,0.004. KD holds T1 retention within
3.3\,pp of the Stage~1 ceiling. Mean cosine similarity S1$\to$S2\,=\,0.917. \\[4pt]

RQ2 &
How does scale shift characterise the forgetting pattern during T3? &
\textbf{Fully answered} &
Large-stratum T1 mAP\,$\to$\,0.000 by epoch~1. Scale distribution
shifts from 85\,\% Normal+Large (Stage~1) to 97.7\,\% Tiny+Small
(Stage~3). Mean cosine similarity S2$\to$S3\,=\,0.968 — forgetting
is gradient starvation, not bulk weight collapse. \\[4pt]

RQ3 &
Does scale-stratified herding mitigate large-target forgetting during T3? &
\textbf{Design contribution} &
300-exemplar buffer built (75 per stratum, iCaRL greedy selection).
Directly addresses gradient starvation diagnosed in RQ2.
Experimental validation is future work. \\
\hline
\end{tabular}
\end{table}

\subsubsection*{H1 — Catastrophic Forgetting Is Measurable (Confirmed)}

\textbf{Hypothesis:} A YOLOMG detector trained sequentially across three
thermal anti-UAV datasets will exhibit measurable catastrophic forgetting
of Task~1 performance.

\textbf{Evidence:}
\begin{itemize}
    \item FM\_naive\,=\,$-$0.605 at peak T3 epoch. T1 mAP collapses
    from 0.640 to 0.035 after 18 epochs of Stage~3 training.
    \item Large-stratum T1 mAP reaches 0.000 by epoch~1 — the model's
    primary detection capability is erased in a single epoch.
    \item For reference, Stage~2 (domain shift only, with KD) produced
    FM\,=\,$-$0.033\,$\pm$\,0.004 — confirming the Stage~3 collapse
    is not a generic property of sequential training but specific to
    the scale-distribution shift introduced by CST Anti-UAV.
\end{itemize}

\noindent \textbf{Status: Confirmed.} Forgetting is measurable, large,
and stratum-specific. This validates the thesis premise and motivates both
RQ2 (characterisation) and RQ3 (mitigation).

\subsubsection*{RQ1 — KD Preserving T1 Performance During T2 Domain Shift (Fully Answered)}

\textbf{Question:} Does knowledge distillation preserve Task~1 performance
during Task~2 training under domain shift? Specifically: does the KD loss
term in Stage~2 (Anti-UAV410) training prevent the model from forgetting
its Anti-UAV-RGBT detection capability, despite the domain shift between
the two datasets?

\textbf{Evidence:}
\begin{itemize}
    \item FM at peak T2 epoch: $-$0.033\,$\pm$\,0.004 across 3 seeds.
    The model retains 95\% of T1 knowledge while adapting to a new
    thermal dataset. KD successfully constrains forgetting.
    \item KD/detection loss ratio rises from 1.36$\times$ (epoch~1) to
    2.60$\times$ (final epoch) — the KD term is active and growing
    throughout training, not a dead loss term.
    \item Parameter drift: mean cosine similarity Stage~1$\to$Stage~2
    = 0.917. The spatial attention module (model.4) shows the highest
    drift ($>4\times$ relative L2), indicating the model adapted its
    attention patterns for the new domain while preserving backbone
    representations — the desired stability-plasticity balance.
\end{itemize}

\noindent \textbf{Status: Fully answered.} FM, KD ratio dynamics, and
parameter drift together provide a quantitative and mechanistic answer.
The original RGB$\to$TIR cross-modal UDA component is repositioned as
future work — Anti-UAV-RGBT's paired RGB streams make it directly
feasible in a follow-up study.

\subsubsection*{RQ2 — Scale Shift Characterising the Forgetting Pattern During T3 (Fully Answered)}

\textbf{Question:} How does scale shift characterise the pattern of
catastrophic forgetting during Task~3 training? Specifically: does the
extreme size distribution of CST Anti-UAV — 67.3\,\% tiny, 30.4\,\%
small, $\approx$0\,\% large — produce stratum-specific forgetting in
which large-target detection collapses disproportionately fast?

\textbf{Evidence:}
\begin{itemize}
    \item \textbf{Scale distribution shift:} Stage~1 (Anti-UAV-RGBT)
    has $\approx$85\,\% Normal+Large targets; Stage~3 (CST Anti-UAV)
    has 97.7\,\% Tiny+Small. The gradient signal for large targets
    is effectively zero during Stage~3 training.
    \item \textbf{Stratum-specific collapse:} Large-stratum T1 mAP
    reaches 0.000 by epoch~1. Normal-stratum collapses shortly after.
    FM\_naive\,=\,$-$0.605 at peak T3 epoch; T1 mAP drops from 0.640
    to 0.035 in 18 epochs.
    \item \textbf{Mechanism — gradient starvation, not weight collapse:}
    Mean cosine similarity S2$\to$S3\,=\,0.968 (higher than the
    S1$\to$S2 shift of 0.917 under KD). The model's weights change less
    in absolute terms than during Stage~2, yet forgetting is 18$\times$
    worse. This confirms the cause is the absence of large-target
    gradient signal, not generalised weight overwriting.
\end{itemize}

\noindent \textbf{Status: Fully answered.} The forgetting pattern is
quantified at the stratum level, the causal mechanism is identified
(gradient starvation from scale distribution mismatch), and the
direction of the solution follows necessarily: re-inject gradient
signal for the starved strata via stratified replay (RQ3).

\subsubsection*{RQ3 — Scale-Stratified Herding Mitigating Large-Target Forgetting (Design Contribution)}

\textbf{Question:} Does scale-stratified herding mitigate large-target
forgetting during Task~3 training? Specifically: does a replay buffer
that selects exemplars proportionally across four UAV size strata
(tiny / small / normal / large) — using iCaRL-style greedy herding
within each stratum — reduce the stratum-specific forgetting identified
in RQ2 compared to naive fine-tuning?

\textbf{What was implemented:}
\begin{itemize}
    \item A 300-exemplar replay buffer with 75 exemplars per stratum,
    selected by greedy iCaRL herding in the Stage~2 embedding space.
    The selection criterion is minimum mean distance to the stratum
    centroid, ensuring each exemplar is maximally representative.
    \item The training loop in \texttt{train\_stage3.py} integrates the
    replay buffer via a separate backward pass at each step (REPLAY\_WEIGHT\,=\,4.0,
    REPLAY\_RATIO\,=\,0.25), verified correct through three DDP crash-fix cycles
    (autograd in-place conflict, logging tensor format, replay image loading bug).
    \item The design follows mechanically from the RQ2 diagnosis: gradient
    starvation of large targets requires re-injecting large-target gradient
    signal, which is exactly what stratum-proportional replay achieves.
\end{itemize}

\textbf{What was not validated:} The herding condition was submitted four times
but cancelled due to budget exhaustion before completing a full training run.
No FM\_herding result exists. The comparison naive vs herding vs random
stratified is future work.

\noindent \textbf{Status: Design contribution.} The mechanism is implemented,
motivated by RQ2, and directly addresses gradient starvation. The thesis
contributes the algorithm and its integration; experimental validation of FM
reduction is the natural next empirical step. The original RQ3 (evidential
uncertainty for tracking) is the logical horizon \textit{after} herding
validation — it is not a dropped idea but a deferred one.

% ------------------------------------------------------------
\subsection*{7. Decisions Made at Supervisor Meeting (5~June~2026)}

\begin{enumerate}
    \item \textbf{RQ structure approved.} The split into H1 + RQ1 + RQ2 + RQ3
    as documented above is the final structure for submission. RQ2 and RQ1 are
    both fully answered. RQ3 is a design contribution with experimental
    validation as future work.

    \item \textbf{Herding framing approved.} RQ3 is presented as a principled
    design contribution — the algorithm, its motivation from RQ2's gradient
    starvation diagnosis, and its implementation. The absence of a completed
    training run is explained honestly; the code is correct.

    \item \textbf{Original RQ3 (evidential uncertainty + IDSW) dropped.}
    Repositioned as the natural future work horizon contingent on herding
    validation. Not mentioned in the results chapter; addressed in the
    future work section only.

    \item \textbf{Results chapter scope:} Present the full three-condition
    design (naive / random stratified / herding) in the methodology chapter.
    Results chapter reports naive only. Herding and random stratified are
    described as unexecuted due to compute budget, with the implementation
    being the contribution.
\end{enumerate}

% ------------------------------------------------------------
\subsection*{8. Timeline to Submission (Revised)}

\begin{center}\small
\begin{tabular}{llp{0.55\linewidth}}
\hline
\textbf{Date} & \textbf{Week} & \textbf{Milestone} \\
\hline
5~June   & Wk 22 & Supervisor meeting — scope sign-off \\
5--9~June & Wk 22--23 & Stage~3 naive analysis; all figures generated \\
9--16~June & Wk 23 & Results chapter (Stage~1, 2, 3 naive) \\
16--22~June & Wk 24 & Methodology, discussion, conclusion; RQ answers \\
22--25~June & Wk 25 & Buffer — polish, proof-read, abstract \\
26~June  & Wk 26 & \textbf{Submission} \\
\hline
\end{tabular}
\end{center}

\noindent No compute milestones remain. All remaining work is analysis
and writing.

% ------------------------------------------------------------
\subsection*{9. Thesis Narrative}

The thesis tells one coherent story, structured as a chain where each
finding motivates the next. The narrative is not ``herding works'' and
not ``here is what I could not do.'' It is:

\begin{quote}
\textit{Scale-distribution shift is the dominant driver of catastrophic
forgetting in sequential thermal anti-UAV detection. We identify the
mechanism, quantify it empirically, and design a principled mitigation
strategy. Robust uncertainty-aware tracking is the natural next chapter,
contingent on the forgetting first being solved.}
\end{quote}

\subsubsection*{9.1 The chain}

\paragraph{Step 1 — Establish that domain adaptation is survivable (RQ1).}
Moving from Anti-UAV-RGBT to Anti-UAV410 with a KD loss produces
FM\,=\,$-$0.033\,$\pm$\,0.004. The model retains 95\% of T1 knowledge.
Parameter drift confirms the mechanism: mean cosine similarity of 0.917,
with spatial attention adapting most ($>4\times$ relative L2) while the
backbone is largely preserved. KD successfully enforces
stability-plasticity balance across a dataset shift. \textit{This
fully answers RQ1 and sets the T1 ceiling against which all forgetting is
measured.}

\paragraph{Step 2 — Show that scale shift is catastrophic and stratum-specific (RQ2).}
Adding CST Anti-UAV (97.7\,\% tiny+small targets, $\approx$0\,\% large)
without any replay produces FM\,=\,$-$0.605 in 18 epochs — 18$\times$
worse than Stage~2. T1 mAP collapses from 0.640 to 0.035. Large-stratum
mAP reaches 0.000 by epoch~1. The mechanism: CST provides essentially
zero gradient signal for large targets, so the model erases those
representations to free capacity for tiny-target detection. Critically,
mean cosine similarity S2$\to$S3\,=\,0.968 — the weights change
\textit{less} than in Stage~2, yet forgetting is 18$\times$ worse.
This confirms \textbf{gradient starvation}, not bulk weight overwriting.
\textit{This fully answers RQ2 and prescribes the solution directly:
re-inject gradient signal for the starved strata.}

\paragraph{Step 3 — Design the direct mechanical response (RQ3).}
Scale-Stratified Herding builds a 300-exemplar buffer (75 per stratum:
tiny / small / normal / large) using greedy iCaRL selection in the Stage~2
feature space. The design follows necessarily from the RQ2 diagnosis — it
is not a generic replay buffer but a direct counter to gradient starvation.
The implementation is complete and verified correct (three DDP crash-fix
cycles resolved). Experimental validation of FM reduction is the natural
next empirical step. \textit{This answers RQ3 at the design level: the
algorithm, its mechanistic motivation, and its integration into the
three-stage curriculum are the contribution.}

\paragraph{Step 4 — Future work: uncertainty-aware tracking.}
Evidential uncertainty for tracking (IDSW reduction during thermal
crossover) is the logical next step \textit{after} herding validation.
Worrying about tracking robustness is premature while the detector is
catastrophically forgetting large targets entirely. Once herding
stabilises T1 performance, the natural open question is whether the
stable detector can handle uncertainty during crossover events.
Anti-UAV-RGBT provides paired RGB streams and crossover annotation flags,
making this directly feasible in a follow-up study. \textit{This is not
a dropped idea — it is the horizon the thesis opens.}

\subsubsection*{9.2 What this means for the writing}

No section of the thesis should mention compute budget, infrastructure
limits, or cancelled experiments. The scope is presented as intentional:
the thesis demonstrates the forgetting problem, diagnoses the mechanism,
and contributes the mitigation design. Validation is a natural next
empirical step. The second examiner sees a complete scientific argument,
not a list of things that did not happen.

% ------------------------------------------------------------
\subsection*{9.3 Remaining Experiments — Everything Needed Before Writing}

Jobs requiring GPU run on a single A100 (128\,SBU/GPU-hour).
CPU-only jobs (training curves, parameter drift, scale distribution)
run on the login node at zero cost. Remaining budget:
$\approx$3,672\,SBU.

\begin{table}[h]
\centering
\caption{Experiments still to run, in priority order.}
\small
\begin{tabular}{p{0.03\linewidth}p{0.32\linewidth}p{0.07\linewidth}p{0.07\linewidth}p{0.40\linewidth}}
\hline
\textbf{\#} & \textbf{Job} & \textbf{Time} & \textbf{SBU} & \textbf{Produces} \\
\hline

1 &
\texttt{eval\_full\_analysis.py} on Stage~3 naive \texttt{best.pt} &
2\,h & 256 &
Scale-stratified mAP (tiny/small/normal/large) on T1 val;
P-R curve; per-sequence dot plot.
\textbf{Ch.4 key table.} \\[4pt]

2 &
\texttt{plot\_training\_analysis.py} on Stage~3 naive \texttt{results.csv} &
$<$1\,min & 0 &
Per-epoch mAP, FM, and per-stratum curves.
Run on login node (CPU only).
\textbf{Ch.4 key figure: stratum collapse. \checkmark Done.} \\[4pt]

3 &
\texttt{parameter\_drift.py} Stage~2 $\to$ Stage~3 naive &
$<$1\,min & 0 &
Per-layer cosine sim and L2 drift; compare to S1$\to$S2.
Run on login node (CPU only).
\textbf{Ch.4: gradient starvation evidence. \checkmark Done.} \\[4pt]

4 &
Scale distribution script across all three datasets &
$<$1\,min & 0 &
Proportion of tiny/small/normal/large per dataset (bar chart).
Run on login node (CPU only).
\textbf{Ch.3: motivates herding design. \checkmark Done.} \\[4pt]

5 &
\texttt{visualise\_detections\_rgbt.py} on Stage~3 naive \texttt{best.pt} &
$<$1\,min & $\approx$2 &
Annotated frames showing T1 detection collapse post-Stage~3.
\textbf{Ch.4: qualitative result. \checkmark Done.} \\[4pt]

6 &
Tracking eval (SR/PR + IDSW) on Stage~3 naive \texttt{best.pt} on CST val sequences &
1.5\,h & 192 &
Success Rate, Precision, IDSW baseline.
\textbf{Ch.5 future work framing. Optional.} \\[4pt]

\hline
\textbf{Total} & & \textbf{$\approx$3.5\,h} & \textbf{$\approx$450} & Well within remaining budget. \\
\hline
\end{tabular}
\end{table}

\noindent \textbf{Note on job 4 (scale distribution):}
This script does not exist yet — it needs to be written. It iterates
over bounding box annotations in each dataset and bins targets by
normalised box area (tiny $<$0.01, small $<$0.04, normal $<$0.1,
large $\geq$0.1 of frame area). Output is one stacked bar chart per
dataset. This is the single most important motivating figure for H1
and RQ2 — it shows visually why CST starves the model of large-target
gradient signal.

\vspace{0.5em}
\noindent \textbf{What each chapter needs:}

\begin{itemize}
    \item \textbf{Ch.3 Methodology:} Scale distribution chart (job~4).
    Herding algorithm description (already written in code, no new job).
    \item \textbf{Ch.4 Results:}
        Stage~1 — scale-stratified mAP table (from job~1, already have
        partial data; confirm on best.pt).
        Stage~2 — FM $\pm$ std, KD ratio, parameter drift S1$\to$S2
        (all already available).
        Stage~3 naive — per-stratum mAP table (job~1), stratum collapse
        figure (job~2), parameter drift S2$\to$S3 (job~3),
        detection visualisations (job~5), FM bar chart (already have numbers).
    \item \textbf{Ch.5 Discussion:} Tracking baseline IDSW (job~6)
    for future work framing. All other numbers already available.
\end{itemize}

\noindent \textbf{After jobs 1--6 are complete, all data needed for
writing exists. No further experiments are required.}

% ------------------------------------------------------------
\subsection*{10. Chapter Writing Plan}

\noindent Total body limit: \textbf{10 pages} (excluding appendix).
All page estimates assume ACM double-column format.

\begin{table}[h]
\centering
\caption{Thesis chapter plan — 10-page budget.}
\small
\begin{tabular}{p{0.04\linewidth}p{0.22\linewidth}p{0.10\linewidth}p{0.54\linewidth}}
\hline
\textbf{Ch.} & \textbf{Title} & \textbf{Pages} & \textbf{Content} \\
\hline
1 & Introduction &
1 &
Problem, stability-plasticity dilemma, gap (no prior work
chains UAV datasets into a forgetting curriculum),
H1 + RQs, one-paragraph contributions summary. \\[4pt]

2 & Related Work &
1.5 &
Four tight paragraphs: continual learning / rehearsal buffers;
UAV detection baselines; domain adaptation (KD);
scale-specific forgetting (gap statement). \\[4pt]

3 & Methodology &
2.5 &
Three-stage curriculum + datasets (1 table);
YOLOMG dual-input architecture (brief);
Stage~2 KD loss;
Scale-Stratified Herding algorithm (1 figure: stratum selection);
FM definition. \\[4pt]

4 & Results &
3.5 &
\textbf{Stage~1} (0.5\,p): T1 ceiling, scale-stratified baseline table.
\textbf{Stage~2} (1\,p): T2 mAP\,=\,0.428\,$\pm$\,0.002,
FM\,=\,$-$0.033\,$\pm$\,0.004, parameter drift figure.
\textbf{Stage~3 naive} (2\,p): FM\,=\,$-$0.605,
per-stratum collapse figure (key result),
H1 comparison (FM\_modality vs FM\_scale bar chart). \\[4pt]

5 & Discussion \& Conclusion &
1.5 &
H1 confirmed (18$\times$ ratio);
RQ1 reframed (KD + drift evidence);
RQ2 partial (forgetting characterised, herding as design contribution);
RQ3 + herding validation as future work;
one-paragraph conclusion. \\
\hline
\textbf{Total} & & \textbf{10\,p} & \\
\hline
\end{tabular}
\end{table}

\noindent \textbf{What goes in the appendix:} full per-epoch result tables,
parameter drift plots for all layers, training loss curves, additional
visualisation frames, dataset statistics tables. Anything that supports
the claims but does not fit the 10-page body.

% ------------------------------------------------------------
\subsection*{11. Analysis and Figure Plan (Remaining Budget)}

Remaining budget: $\approx$3,672\,SBU. A100 partition:
128\,SBU/GPU-hour (single GPU). CPU-only jobs run on the login
node at zero cost.

\begin{table}[h]
\centering
\caption{Analysis jobs — actual status.}
\small
\begin{tabular}{p{0.38\linewidth}p{0.08\linewidth}p{0.08\linewidth}p{0.36\linewidth}}
\hline
\textbf{Job} & \textbf{Time} & \textbf{SBU} & \textbf{Output} \\
\hline
\texttt{eval\_full\_analysis.py} on Stage~3 naive \texttt{best.pt}
& 2\,h & 256 &
Scale-stratified mAP (tiny/small/normal/large), P-R curve,
per-sequence dot plot. \textit{Running (job 23821658).} \\[4pt]

\texttt{plot\_training\_analysis.py} on Stage~3 naive \texttt{results.csv}
& $<$1\,min & 0 &
mAP/loss/FM per epoch curves; stratum collapse figure.
\textit{\checkmark Done (login node).} \\[4pt]

\texttt{parameter\_drift.py} Stage~2 $\to$ Stage~3 naive
& $<$1\,min & 0 &
Per-layer cosine sim\,=\,0.968; compare to S1$\to$S2 (0.917).
\textit{\checkmark Done (login node).} \\[4pt]

\texttt{plot\_scale\_distribution.py} across all three datasets
& $<$1\,min & 0 &
Scale distribution bar chart.
\textit{\checkmark Done (login node).} \\[4pt]

\texttt{visualise\_detections\_rgbt.py} on Stage~3 naive \texttt{best.pt}
& 26\,s & $\approx$2 &
300 annotated frames showing detection collapse on T1 sequences.
\textit{\checkmark Done (job 23821661).} \\[4pt]
\hline
\textbf{Total} & & \textbf{$\approx$258} & Well within budget. \\
\hline
\end{tabular}
\end{table}

\noindent \textbf{Figures for the thesis} (generated from the above):

\begin{enumerate}
    \item FM comparison bar chart: Stage~2 vs Stage~3 naive
    ($-$0.033 vs $-$0.605), with error bars for Stage~2 (n=3 seeds).
    \item T1 mAP per epoch (Stage~3 naive) with four stratum curves
    overlaid — shows large-target erasure by epoch~2.
    \item Scale distribution figure: three datasets side-by-side
    (tiny/small/normal/large proportions). Motivates H1.
    \item Parameter drift plots: S1$\to$S2 vs S2$\to$S3 — shows
    scale shift moves weights further than modality shift.
    \item Training loss curves: Stage~1, 2, 3 (CST loss + replay loss).
    \item PR curves: Stage~1 and Stage~2 best.pt on T1 val.
\end{enumerate}

% ------------------------------------------------------------
\subsection*{12. RQ3 Implementation Analysis — Future Work}

This section provides a detailed breakdown of what RQ3 would require
to implement, for the supervisor's information. The goal is to make
the scope decision concrete rather than abstract.

\subsubsection*{12.1 What RQ3 Requires}

RQ3 asked: \textit{``How does the integration of Evidential Uncertainty
and Bimodal Fusion reduce identity switches (IDSW) during thermal
crossover events?''}

Fully answering RQ3 requires three independent components:

\begin{enumerate}
    \item \textbf{Real motion signal in \texttt{img2}} — currently
    \texttt{img2\,=\,zeros} throughout. To activate the motion branch,
    inter-frame motion difference maps must be precomputed for each
    training dataset and the dataset loaders modified to load them.

    \item \textbf{Evidential Head} — replaces the detection head's
    classification output with Dirichlet distribution parameters. Adds
    a KL divergence regularisation term to the loss. Computes vacuity
    (epistemic uncertainty) per frame.

    \item \textbf{Uncertainty-Weighted Fusion Gate} — uses vacuity to
    dynamically down-weight the appearance branch during thermal crossover.
    Requires the two input streams to remain \textit{separate} until late
    in the network so they can be independently gated. Currently YOLOMG
    merges both inputs in the first convolutional layer, making late-stage
    gating architecturally impossible without refactoring.
\end{enumerate}

\noindent \textbf{All three components are required together.}
Implementing any one or two in isolation does not answer RQ3:
\begin{itemize}
    \item Motion masks without evidential head = standard YOLOMG bimodal
    fusion (fixed-weight, no crossover awareness). Answers ``does motion
    help detection?'' not ``does uncertainty gating reduce IDSW?''
    \item Evidential head without motion = uncertainty quantification on
    appearance-only detection. Useful, but the vacuity has nothing to
    gate to.
    \item Fusion gate without evidential head = arbitrary, non-principled
    weighting.
\end{itemize}

\subsubsection*{12.2 Implementation Difficulty}

\begin{table}[h]
\centering
\caption{RQ3 component implementation difficulty and time estimate.}
\label{tab:rq3_effort}
\small
\begin{tabular}{p{0.28\linewidth}p{0.14\linewidth}p{0.50\linewidth}}
\hline
\textbf{Component} & \textbf{Effort} & \textbf{Notes} \\
\hline
Motion mask extraction (Anti-UAV-RGBT)
& 1 day
& \texttt{generate\_masks\_npz.py} already exists. Run ECC + frame
  diff over all sequences. Anti-UAV-410 blocked (images only).
  CST trivial (static camera, simple diff). \\[4pt]
Dataset loader modifications
& 0.5 days
& Load precomputed \texttt{img2} from npz instead of zeros.
  Straightforward code change. \\[4pt]
Evidential head (Dirichlet output + KL loss)
& 2--3 days
& Well-documented in Sensoy et al.\ (2018). Add to YOLOMG detection
  head. Modify loss function. Medium difficulty. \\[4pt]
Uncertainty-weighted fusion gate
& 5--7 days
& \textbf{Hard.} YOLOMG fuses both input channels in layer~0
  (concatenation). Separating the streams to enable late-stage
  gating requires refactoring the backbone. High risk of introducing
  bugs or degrading base performance. \\[4pt]
Retrain S1 + S2 + S3
& $\approx$3.5 days compute
& S1 $\sim$14\,h, S2 $\sim$30\,h (1 seed), S3 $\sim$30\,h $\times$3 \\[4pt]
IDSW evaluation on crossover sequences
& 1 day
& Write tracking evaluation script (IoU-Hungarian matching). \\
\hline
\textbf{Total (full RQ3)} & \textbf{$\sim$2 weeks} & Inclusive of compute time, debugging, and validation. \\
\hline
\end{tabular}
\end{table}

\subsubsection*{12.3 Partial Implementation Options}

If full RQ3 is not feasible, two partial options exist:

\paragraph{Option A — Motion masks only (no evidential head).}
Activate \texttt{img2} with real motion maps for Anti-UAV-RGBT and CST.
Retrain all three stages. Report whether bimodal input improves aggregate
mAP and per-stratum detection. \textbf{Does not answer RQ3} (no crossover
gating, no IDSW measurement). Answers a different, weaker question:
``does motion information improve UAV detection?''
\begin{itemize}
    \item Effort: $\sim$1 week (masks + loaders + retrain)
    \item Anti-UAV-410 remains \texttt{img2\,=\,zeros} (images only — unavoidable)
    \item Pipeline inconsistency: motion active in S1 and S3 but not S2
\end{itemize}

\paragraph{Option B — Evidential head only (no motion, no fusion gate).}
Add Dirichlet output and KL loss to the existing appearance-only model.
Provides uncertainty calibration (ECE metric). Vacuity on crossover
sequences could be measured as a proxy for crossover detection.
\textbf{Partially addresses RQ3} — shows the model ``knows when it
doesn't know'' but does not demonstrate performance improvement.
\begin{itemize}
    \item Effort: $\sim$3--4 days (coding + retrain S3 only, no S1/S2 changes)
    \item No motion masks needed — img2 stays zeros
    \item Adds ECE calibration metric to the evaluation
\end{itemize}

\subsubsection*{12.4 Recommendation and Decision Required}

\begin{table}[h]
\centering
\caption{RQ3 implementation options vs.\